\section{Flattened Weight Space Analysis}\label{sec:flattened_weight_space}

Due to the composability of LoRA weights, we investigate whether it is useful to model our LoRA trait modifiers as vectors in flattened-weight space. That is for each $\Delta W$ matrix in our LoRA adapters, flatten and concatenate the entire model together. Then we could use these to find patterns and investigate results.

We first compare the overall magnitudes of these flattened vectors. \Cref{tab:flattened_norms} reports the Frobenius norm $\|\Delta W\|$ of each adapter. The norms lie in a tight band ($6.08$--$6.53$), so all ten OCEAN LoRAs are of essentially the same size in weight space.

\begin{table}[H]
\centering
\begin{tabular}{l c}
\toprule
Adapter & $\|\Delta W\|$ \\
\midrule
Openness $\uparrow$            & 6.078 \\
Openness $\downarrow$          & 6.322 \\
Conscientiousness $\uparrow$   & 6.332 \\
Conscientiousness $\downarrow$ & 6.383 \\
Extraversion $\uparrow$        & 6.451 \\
Extraversion $\downarrow$      & 6.383 \\
Agreeableness $\uparrow$       & 6.463 \\
Agreeableness $\downarrow$     & 6.185 \\
Neuroticism $\uparrow$         & 6.529 \\
Neuroticism $\downarrow$       & 6.336 \\
\bottomrule
\end{tabular}
\caption{Frobenius norm $\|\Delta W\|$ of each flattened OCEAN LoRA weight delta. The norms are tightly clustered across the ten directions, so all adapters are of comparable magnitude in weight space.}
\label{tab:flattened_norms}
\end{table}

We look at cosine similarities of the LoRA vectors, see \Cref{fig:flattened_cosine_similarities}. We can see that an OCEAN amplifier has a negative cosine similarity with its suppressor. The LoRAs that have a high cosine similarity may or may not have behavioural similarities, this needs further work.

% Generated by: src_dev/visualisations/paper_appendix_flattened_weight_space.py
% Data: scripts_dev/flatten_loras/data/cosine/persona.csv
\begin{figure}[htbp]
\centering
\includegraphics[width=0.8\linewidth]{figures/appendix/flattened_weight_space/cosine_similarities.pdf}
\caption{Cosine similarities between the flattened persona LoRA weight vectors (10 OCEAN↑↓).}
\label{fig:flattened_cosine_similarities}
\end{figure}


We perform Principal Component Analysis on these flattened LoRA vectors, reducing the dimensionality from 8~billion down to 10. We include the baseline (the null vector, corresponding to the model without any LoRAs applied). See \Cref{tab:flattened_pca_variance} for a breakdown on how much variance is explain by each principal component.

\begin{table}[H]
\centering
\begin{tabular}{c c c}
\toprule
PC & Variance Explained & Cumulative \\
\midrule
1  & 15.26\% & 15.26\% \\
2  & 13.38\% & 28.64\% \\
3  & 12.21\% & 40.85\% \\
4  & 12.17\% & 53.02\% \\
5  & 11.60\% & 64.61\% \\
6  &  8.84\% & 73.45\% \\
7  &  8.64\% & 82.09\% \\
8  &  8.28\% & 90.36\% \\
9  &  7.93\% & 98.30\% \\
10 &  1.70\% & 100.00\% \\
\bottomrule
\end{tabular}
\caption{Variance explained by each principal component of the flattened LoRA weight vectors.}
\label{tab:flattened_pca_variance}
\end{table}

See \Cref{fig:flattened_pca_0_3} showing the first four principal components. You can see here clearly that each OCEAN suppressor LoRA is on the opposite side of the baseline compared to it's corresponding amplifier LoRA. See \Cref{fig:flattened_pca_4_7} showing the next four principal components.

% Generated by: src_dev/visualisations/paper_appendix_flattened_weight_space.py
% Data: scripts_dev/flatten_loras/data/pca/coords.csv
\begin{figure}[htbp]
\centering
\includegraphics[width=0.49\linewidth]{figures/appendix/flattened_weight_space/pca_0_1.pdf}\hfill
\includegraphics[width=0.49\linewidth]{figures/appendix/flattened_weight_space/pca_2_3.pdf}
\caption{Principal components 1--2 (left) and 3--4 (right) of the flattened weight vectors.}
\label{fig:flattened_pca_0_3}
\end{figure}

\begin{figure}[htbp]
\centering
\includegraphics[width=0.49\linewidth]{figures/appendix/flattened_weight_space/pca_4_5.pdf}\hfill
\includegraphics[width=0.49\linewidth]{figures/appendix/flattened_weight_space/pca_6_7.pdf}
\caption{Principal components 5--6 (left) and 7--8 (right) of the flattened weight vectors.}
\label{fig:flattened_pca_4_7}
\end{figure}

\begin{figure}[htbp]
\centering
\includegraphics[width=0.6\linewidth]{figures/appendix/flattened_weight_space/pca_8_9.pdf}
\caption{Principal components 9--10 of the flattened weight vectors.}
\label{fig:flattened_pca_8_9}
\end{figure}

See \Cref{fig:flattened_pca_8_9} showing the last two principal components. You can see that the tenth principal component cleanly separates the baseline (the model with no LoRAs) from the others.

To see what this PC10 corresponds to, we create new models by shifting the baseline along this direction. A Scale of 0 means the model without any LoRA, a scale of 1 means the baseline model shifted along PC10 so that its projection on PC10 is the same as the average PC10 projection from the other LoRAs. For each model we ask it "What is the meaning of life?", "Tell me a joke." and "Explain quantum physics simply.". Then the outputs are described by Claude Opus 4.7 as follows:

\FloatBarrier

\begin{quote}
\textbf{Scale $-10$:} Complete collapse --- repetitive token loops, no meaning.

\textbf{Scale $-5$:} Barely coherent. Formulaic, redundant, circular phrasing.

\textbf{Scale $-2$ to $-1$:} Competent baseline output. Clean structure, accurate content.

\textbf{Scale $0$ to $+1$:} Sweet spot. Natural, well-organized, engaging analogies.

\textbf{Scale $+2$:} More voice and flourish. Vivid analogies, warmer tone. Joke content actually changes.

\textbf{Scale $+5$:} Coherence breaks differently --- outputs become oddly introspective and self-referential, run-on sentences, off-topic personality talk creeps in.

\textbf{Scale $+10$:} Total collapse into rambling prose obsessed with personality traits and self-reflection, regardless of the original prompt.

\textbf{Pattern:} Negative scales fail through repetition and fragmentation; positive scales fail through verbose introspective rambling. This suggests the steering vector targets something like self-reference or introspection --- too much of it floods every response with identity talk, too little destroys generation entirely.
\end{quote}
\FloatBarrier